\scanpage{17}
\margin{faire le raisonnement}% slanted

Alors on voit que $P_{\mathrm{an}}(1_Y)$ est la classe de \unsure{Ch.|Coh.} sur $X-a$ représentée par la forme diff. $\frac{dz}{z}$, et par suite \deleted{\unsure{on a}} si $P_{\mathrm{an}}(1_Y)=\lambda P_{\mathrm{top}}(1_Y)$, on a
\[
\lambda=\underset{\substack{\nwarrow\qquad\nearrow\\\text{\inserted{cercle \unsure{centré|anti} \deleted{\unsure{cercle|en $a$}} orienté.}}}}{\int_C P_{\mathrm{an}}(1_Y)=\int_C\frac{dz}{z}}=2i\pi
\]

Donc
\[
\boxed{P_{\mathrm{an}}\left(\frac{1}{2i\pi}1_Y\right)=P_{\mathrm{top}}(1_Y)}
\]
\[
\boxed{\left\{\begin{aligned}
P_{\mathrm{top}}(1_Y)&=\frac{1}{2i\pi}P_{\mathrm{an}}(1_Y)\\
P_{\mathrm{an}}(1_Y)&=2i\pi\,P_{\mathrm{top}}(1_Y)
\end{aligned}\right.}
\]

\uline{\unsure{Cor.|Gys|Gén.|Ex.}} Considérons le diagramme
\diagram{\resizebox{\linewidth}{!}{\begin{tikzcd}[ampersand replacement=\&,row sep=large,column sep=large]
H_{\unsure{\alpha|a}}^{i+p}(X,\mathbb Z) \arrow[r,leftrightarrow] \& \unsure{H|\mathbb H}_{\deleted{\unsure{Z|Y}}}^{\deleted{\unsure{2|j}}i+p}(X,\mathbb C) \arrow[r,dash,"{\substack{\text{De Rham}\\\sim}}"] \& H^{\deleted{\unsure{2|j}}i+p}(X,\Omega_X^\bullet)\\
H_{\deleted{\unsure{Z|Y}}}^{i}(Y,\mathbb Z) \arrow[u,"{P_{\mathrm{top}}}"] \arrow[r] \& H_{\deleted{\unsure{Z|Y}}}^{\deleted{\unsure{2|j}}i}(\deleted{\unsure{X}}Y,\mathbb C) \arrow[u,"{P_{\mathrm{top}}}"] \arrow[r,dash,"{\substack{\deleted{\unsure{\text{De R.}}}\\\sim}}","{\frac{1}{(2i\pi)^p}\,\mathrm{can.}}"'] \& \underline{\underline H}_{\deleted{\unsure{Z|Y}}}^{i}(Y,\Omega^\bullet_Y) \arrow[u,"{P_{\mathrm{an}}}"]
\end{tikzcd}}}

\uline{Il est commutatif}, i.e. la formule ci-dessus est valable \uline{globalement}, et en remplaçant $1_Y$ par une classe de Coh. quelc. Provient de la comm. des
% archive sheet number: 16
